%% PLANTILLA PARA ARTICULOS

\documentclass[10pt,twoside,spanish,a4paper]{article}
%%%%%%%%%%%%%%%%%%%%%%%%%%%%% no tocar
% Alfredo Cuesta. 01/01/2013. acuestai@ucm.es

%si se quiere usar ARIAL hay que poner estas 2 lineas:
%\usepackage[scaled]{uarial} %DESCOMENTAR
\renewcommand*\familydefault{\sfdefault}
%

\usepackage[compact]{titlesec}
 \titlespacing\section{0pt}{12pt plus 0pt minus 0pt}{0pt plus 0pt minus 0pt}
 \titlespacing\subsection{0pt}{5pt plus 0pt minus 0pt}{0pt plus 0pt minus 0pt}
 \titlespacing\subsubsection{0pt}{0pt plus 0pt minus 0pt}{0pt plus 0pt minus 0pt}


%\usepackage[T1]{fontenc}
\usepackage[utf8]{inputenc}
%\usepackage[latin9]{inputenc}
\usepackage{
amssymb,
graphicx,
color,
url
%wrapfig %Inclusión de gr·ficos al lado de texto
}

\linespread{1.3}%separacion de lineas
\parindent=0.0in%indentancion de parrafo
\parskip=6pt%separacion entre el parrafo siguiente
\setlength{\topmargin}{-0.5in}
\setlength{\oddsidemargin}{0in}
\setlength{\evensidemargin}{0in}
\setlength{\textwidth}{16cm}%{4.72in}
\setlength{\textheight}{22cm}%{7.08in}

\usepackage{babel}
\addto\shorthandsspanish{\spanishdeactivate{~<>}}
% correct bad hyphenation here
\hyphenation{nu-es-tro co-no-ci-mien-to par-ti-cu-lar ca-rac-te-ris-ti-cas ni-vel se-gui-mien-to que-re-mos ca-ta-lo-go}
\date{}


% --- customization ---

\renewcommand\labelitemi{$\bullet$}
\let\olditemize\itemize
\renewcommand{\itemize}{
  \olditemize
  \vspace{-0pt}
  \setlength{\itemsep}{6pt}
  \setlength{\parskip}{0pt}
  \setlength{\parsep}{0pt}
  %\setlength{\topsep}{-15pt}
}

\let\oldenumerate\enumerate
\renewcommand{\enumerate}{
  \oldenumerate
  \vspace{-2pt}
  \setlength{\itemsep}{1pt}
  \setlength{\parskip}{0pt}
  \setlength{\parsep}{0pt}
  %\setlength{\topsep}{-15pt}
}

\usepackage{verbatim}
\usepackage{fancyhdr}
\usepackage{multirow}
\usepackage{colortbl}
\usepackage{verbatim}

\definecolor{olivegreen}{rgb}{0.0,0.5,0.0}

\begin{document}

\renewcommand{\tablename}{Tabla}
\renewcommand{\figurename}{Figura}
\title{
  \vspace{-0.2cm}
	\begin{center}
		\large
		\begin{tabular}{rrl}
		\multirow{4}{*}{\includegraphics[width=1.7cm]{UCM.eps} \hspace{1cm}}&Grado: & Ingeniería Electrónica de Comunicaciones\\
		&Asignatura:& Procesamiento de Señales \\
		&Profesores:  & Jesús Chacón Sombría \\
		&             & María José Gómez Silva \\
		\end{tabular}
	\end{center}
	\rule{\textwidth}{0.4pt}
	%
	\vspace{-1.5cm}
	\begin{center}
		\color{blue} \Huge
		\begin{tabular}{c}
	        \textbf{Ejercicios}\\
		\textbf{Muestreo en el dominio de la frecuencia}\\
		\end{tabular}
	\end{center}
	\vspace{-0.5cm}
	\rule{\textwidth}{0.4pt}
}
\maketitle
 \thispagestyle{empty}
\pagestyle{fancy}

\lfoot{\includegraphics[width=0.8cm]{UCM.eps}}
\cfoot{\vspace{-0.6cm} Ingeniería Electrónica de Comunicaciones }
\rfoot{\vspace{-0.6cm} \thepage}
\lhead{Ejercicios}
\rhead{Muestreo en el dominio del tiempo}
\renewcommand{\headrule}{ \vspace{-0.2cm} \rule{\textwidth}{0.4pt}}
\renewcommand{\footrule}{\vspace{0.5cm} \rule{\textwidth}{0.4pt}\\ \vspace{-0.3cm} \\}

%**************************************************************<- para indicar pregunta de la pr·ctica
\newcounter{ctPreg}
\setcounter{ctPreg}{1}
\newcommand{\itemP}{ \item[{\color{blue}\fbox{\textbf{Ejercicio \arabic{ctPreg}:}}}] \addtocounter{ctPreg}{1}}  %
\newcommand{\itemB}{ \item[{\color{blue}$\bullet$} ]    }
\newcommand{\itemR}{ \item[{\color{blue}$\triangleright$} ]    }       %
%**************************************************************

\vspace{-8em}

	\begin{itemize}
		\itemP Para cada señal, calcula y representa su DFT de $N$ muestras en MATLAB (mediante \color{blue}\texttt{fft}\color{black}). Indica el valor de $N$ y del periodo de muestreo $T_s$ ($t=n T_s$) y justifica su elección.
			\begin{itemize}
				\itemB $x(t)=\textrm{sin}(\frac{\pi}{2}t)$
				\itemB $x(t)=\sum_{k=0}^{10}0.9^k sin(\frac{k\pi}{2}t)$
				\itemB $x_1(t)=\textrm{e}^{-t}, 0\leq t<5$
				\itemB $x_2(t)=\textrm{sinc}(t)$
				\itemB $x_2(t)=\textrm{randn}(N,1)$
			\end{itemize}

		\itemP Reconstruye las señales del ejercicio 1 a partir de la DFT, en los casos en que sea posible.

		\itemP Reconstruye la DTFT de las señales del ejercicio 1 a partir de la DFT, en los casos en que sea posible.
		
		\itemP Reconstruye la transformada Z de las señales del ejercicio 1 a partir de la DFT, en los casos en que sea posible.

		\itemP Calcula, mediante la DFT, la respuesta de los siguientes sistemas LTI tomando como entrada las señales del ejercicio 1. 
		\begin{itemize}
			\itemB $h(n)=\frac{1}{N}, 0\leq n < N=20$
			\itemB $h(n)=0.7^{-n}$
			\itemB $H(z)=0.8 + 0.5z^{-1} - 0.2 z^{-2} - 4 z^{-3} + z^{-4}$
			\itemB $H(z)=\frac{0.8}{1 + 0.5z^{-1}}$
		\end{itemize}
	
\end{document}
